% Môn: Vật lý
% Lớp: 12
% Chương: Chương I. Vật lí nhiệt
% Bài: L12C1 Bài 4. Nhiệt dung riêng · Khái niệm, công thức, đơn vị và ý nghĩa của nhiệt hóa hơi riêng
% ===== Câu 61 =====
% ID: L12C1B4-115-DS
% Mức: NB
\dangbt{Khái niệm, công thức, đơn vị và ý nghĩa của nhiệt hóa hơi riêng}
\begin{ex}
% ID: L12C1B4-115-DS
% Mức: NB
Xét sự sôi của chất lỏng tinh khiết ở áp suất không đổi.
\choiceTF
{\True Sự hóa hơi xảy ra cả trong lòng và trên mặt thoáng.}
{\True Nhiệt lượng vẫn được chất nhận trong khi nhiệt độ gần như không đổi.}
{Mọi chất lỏng đều có cùng nhiệt độ sôi.}
{\True Nhiệt hóa hơi riêng đặc trưng cho năng lượng chuyển thể trên một đơn vị khối lượng.}
\loigiai{
\begin{itemchoice}
\item (Đúng) Sự hóa hơi xảy ra cả trong lòng và trên mặt thoáng.
\item (Đúng) Nhiệt lượng vẫn được chất nhận trong khi nhiệt độ gần như không đổi.
\item (Sai) Mọi chất lỏng đều có cùng nhiệt độ sôi.
\item (Đúng) Nhiệt hóa hơi riêng đặc trưng cho năng lượng chuyển thể trên một đơn vị khối lượng.
\end{itemchoice}
}
\end{ex}

% ===== Câu 65 =====
% ID: L12C1B4-116-DS
% Mức: TH
\begin{ex}
% ID: L12C1B4-116-DS
% Mức: TH
So sánh hóa hơi và ngưng tụ ở cùng điều kiện.
\choiceTF
{\True Hai quá trình ngược chiều nhau.}
{\True Hóa hơi thu nhiệt, ngưng tụ tỏa nhiệt.}
{\True Độ lớn nhiệt lượng cho cùng khối lượng là như nhau nếu bỏ qua tổn thất.}
{Cả hai quá trình đều làm nhiệt độ tăng liên tục.}
\loigiai{
\begin{itemchoice}
\item (Đúng) Hai quá trình ngược chiều nhau.
\item (Đúng) Hóa hơi thu nhiệt, ngưng tụ tỏa nhiệt. (Vì): a phát biểu đầu đúng; trong chuyển thể ở áp suất không đổi, nhiệt độ gần như không đổi
\item (Đúng) Độ lớn nhiệt lượng cho cùng khối lượng là như nhau nếu bỏ qua tổn thất.
\item (Sai) Cả hai quá trình đều làm nhiệt độ tăng liên tục.
\end{itemchoice}
}
\end{ex}

% ===== Câu 69 =====
% ID: L12C1B4-117-DS
% Mức: VD
\begin{ex}
% ID: L12C1B4-117-DS
% Mức: VD
Về nhiệt hóa hơi riêng của nước, hãy xác định đúng hoặc sai.
\choiceTF
{\True $L$ có đơn vị J/kg.}
{\True Q=Lm dùng cho phần hóa hơi ở nhiệt độ sôi.}
{Trong khi sôi, nhiệt độ nước tăng liên tục nếu tiếp tục cấp nhiệt.}
{\True Ngưng tụ cùng khối lượng hơi ở cùng điều kiện tỏa nhiệt có độ lớn Lm.}
\loigiai{
\begin{itemchoice}
\item (Đúng) $L$ có đơn vị J/kg. (Vì): theo đơn vị SI; b đúng theo định nghĩa; c sai vì nhiệt độ gần như không đổi khi sôi; d đúng theo tính thuận nghịch của chuyển thể
\item (Đúng) Q=Lm dùng cho phần hóa hơi ở nhiệt độ sôi.
\item (Sai) Trong khi sôi, nhiệt độ nước tăng liên tục nếu tiếp tục cấp nhiệt.
\item (Đúng) Ngưng tụ cùng khối lượng hơi ở cùng điều kiện tỏa nhiệt có độ lớn Lm.
\end{itemchoice}
}
\end{ex}

% ===== Câu 73 =====
% ID: L12C1B4-118-DS
% Mức: VD
\begin{ex}
% ID: L12C1B4-118-DS
% Mức: VD
Xét công thức Q=Lm.
\choiceTF
{\True $Q$ tính bằng $J$ nếu $L$ tính bằng J/kg và m tính bằng kg.}
{Có thể dùng m theo gam mà không đổi đơn vị.}
{\True $Q$ tỉ lệ thuận với m khi $L$ không đổi.}
{Công thức này bao gồm cả nhiệt lượng đun từ nhiệt độ ban đầu đến nhiệt độ sôi.}
\loigiai{
\begin{itemchoice}
\item (Đúng) $Q$ tính bằng $J$ nếu $L$ tính bằng J/kg và m tính bằng kg.
\item (Sai) Có thể dùng m theo gam mà không đổi đơn vị.
\item (Đúng) $Q$ tỉ lệ thuận với m khi $L$ không đổi.
\item (Sai) Công thức này bao gồm cả nhiệt lượng đun từ nhiệt độ ban đầu đến nhiệt độ sôi.
\end{itemchoice}
}
\end{ex}

% ===== Câu 74 =====
% ID: L12C1B4-119-TL
% Mức: NB
\begin{bt}
% ID: L12C1B4-119-TL
% Mức: NB
Chứng minh từ công thức Q=Lm rằng với cùng chất và cùng điều kiện, $Q$ tỉ lệ thuận với m. Nêu một hệ quả thực tế.
\loigiai{
Vì $L$ không đổi trong cùng điều kiện nên Q/m= $L$ là hằng số, suy ra $Q$ tỉ lệ thuận m. Do đó muốn đun bay hơi lượng nước gấp đôi cần năng lượng xấp xỉ gấp đôi nếu các điều kiện và tổn thất như nhau.
}
\end{bt}

% ===== Câu 75 =====
% ID: L12C1B4-120-TLN
% Mức: NB
\begin{ex}
% ID: L12C1B4-120-TLN
% Mức: NB
Hóa hơi $0,40\,\text{kg}$ chất lỏng cần 320 kJ. Tính $L$ theo đơn vị 10^ $5\,\text{J}$ /kg. Chỉ ghi số.
\shortans{8}
\loigiai{
$L=Q/m=320000/0,40=8,0.10^5 J/kg.$
}
\end{ex}

% ===== Câu 77 =====
% ID: L12C1B4-121-DS
% Mức: VDC
\begin{ex}
% ID: L12C1B4-121-DS
% Mức: VDC
Một chất lỏng có L=8,0.10^ $5\,\text{J}$ /kg. Xét các phát biểu.
\choiceTF
{\True Hóa hơi $1\,\text{kg}$ ở nhiệt độ sôi cần 8,0.10^ $5\,\text{J}$.}
{\True Hóa hơi $0,5\,\text{kg}$ cần 4,0.10^ $5\,\text{J}$.}
{$L$ là nhiệt dung riêng của chất.}
{\True Ngưng tụ $2\,\text{kg}$ hơi tỏa 1,6.10^ $6\,\text{J}$.}
\loigiai{
\begin{itemchoice}
\item (Đúng) Hóa hơi $1\,\text{kg}$ ở nhiệt độ sôi cần 8,0.10^ $5\,\text{J}$.
\item (Đúng) Hóa hơi $0,5\,\text{kg}$ cần 4,0.10^ $5\,\text{J}$.
\item (Sai) $L$ là nhiệt dung riêng của chất.
\item (Đúng) Ngưng tụ $2\,\text{kg}$ hơi tỏa 1,6.10^ $6\,\text{J}$. (Vì): ùng Q=Lm. $L$ không phải nhiệt dung riêng; hai đại lượng có đơn vị và ý nghĩa khác nhau
\end{itemchoice}
}
\end{ex}

% ===== Câu 78 =====
% ID: L12C1B4-122-TL
% Mức: TH
\begin{bt}
% ID: L12C1B4-122-TL
% Mức: TH
Nêu định nghĩa, kí hiệu, đơn vị SI của nhiệt hóa hơi riêng và giải thích ý nghĩa của L=2,3.10^ $6\,\text{J}$ /kg đối với nước.
\loigiai{
Nhiệt hóa hơi riêng $L$ là nhiệt lượng cần cung cấp để hóa hơi hoàn toàn $1\,\text{kg}$ chất lỏng ở nhiệt độ sôi. Đơn vị SI là J/kg. L=2,3.10^ $6\,\text{J}$ /kg nghĩa là cần 2,3.10^6 $J$ để hóa hơi hoàn toàn $1\,\text{kg}$ nước ở nhiệt độ sôi trong điều kiện xét.
}
\end{bt}

% ===== Câu 79 =====
% ID: L12C1B4-123-TLN
% Mức: NB
\begin{ex}
% ID: L12C1B4-123-TLN
% Mức: NB
Nếu khối lượng hóa hơi tăng từ $0,20\,\text{kg}$ lên $0,50\,\text{kg}$ với $L$ không đổi, nhiệt lượng tăng bao nhiêu lần? Chỉ ghi số.
\shortans{02/05/2026}
\loigiai{
Do $Q$ tỉ lệ với m, Q2/Q1=0,50/0,20=2,5.
}
\end{ex}

% ===== Câu 82 =====
% ID: L12C1B4-124-TL
% Mức: TH
\begin{bt}
% ID: L12C1B4-124-TL
% Mức: TH
Một học sinh nói: “Q=Lm dùng cho mọi giai đoạn đun nước”. Hãy nhận xét và sửa lại.
\loigiai{
Nhận định sai. Q=Lm chỉ tính nhiệt lượng cho giai đoạn nước ở nhiệt độ sôi chuyển thành hơi. Nếu nước ban đầu dưới nhiệt độ sôi, phải cộng thêm Q1=mc $(t_s-t_0)$; tổng Q=mc $(t_s-t_0)$ +Lm.
}
\end{bt}

% ===== Câu 83 =====
% ID: L12C1B4-125-TLN
% Mức: TH
\begin{ex}
% ID: L12C1B4-125-TLN
% Mức: TH
Ngưng tụ $0,20\,\text{kg}$ hơi có L=2,4.10^ $6\,\text{J}$ /kg. Nhiệt lượng tỏa ra theo kJ là bao nhiêu? Chỉ ghi số.
\shortans{480}
\loigiai{
$Q=Lm=2,4.10^6.0,20=480000 J=480 kJ.$
}
\end{ex}

% ===== Câu 86 =====
% ID: L12C1B4-126-TL
% Mức: VD
\begin{bt}
% ID: L12C1B4-126-TL
% Mức: VD
Phân biệt nhiệt dung riêng c và nhiệt hóa hơi riêng $L$ về ý nghĩa, đơn vị và công thức sử dụng.
\loigiai{
c đặc trưng nhiệt lượng làm $1\,\text{kg}$ chất tăng 1 $K$, đơn vị J/(kg.K), dùng Q=mcΔt. $L$ đặc trưng nhiệt lượng hóa hơi hoàn toàn $1\,\text{kg}$ chất lỏng ở nhiệt độ sôi, đơn vị J/kg, dùng Q=Lm.
}
\end{bt}

% ===== Câu 87 =====
% ID: L12C1B4-127-TLN
% Mức: VD
\begin{ex}
% ID: L12C1B4-127-TLN
% Mức: VD
Cần 690 kJ để hóa hơi hoàn toàn chất lỏng có L=2,3.10^ $6\,\text{J}$ /kg. Khối lượng theo gam là bao nhiêu? Chỉ ghi số.
\shortans{300}
\loigiai{
$m=Q/L=690000/(2,3.10^6)=0,30 kg=300 g.$
}
\end{ex}

% ===== Câu 90 =====
% ID: L12C1B4-128-TL
% Mức: VD
\begin{bt}
% ID: L12C1B4-128-TL
% Mức: VD
Giải thích mối liên hệ năng lượng giữa hóa hơi và ngưng tụ của cùng một khối lượng chất ở cùng điều kiện.
\loigiai{
Hóa hơi cần nhận Q=Lm để các phân tử tách xa nhau. Khi ngưng tụ, hệ chuyển ngược lại và tỏa ra môi trường nhiệt lượng có cùng độ lớn Lm nếu trạng thái đầu-cuối tương ứng và bỏ qua các phần trao đổi khác.
}
\end{bt}

% ===== Câu 91 =====
% ID: L12C1B4-129-TLN
% Mức: VDC
\begin{ex}
% ID: L12C1B4-129-TLN
% Mức: VDC
Một chất lỏng có L=2,0.10^ $6\,\text{J}$ /kg. Hóa hơi hoàn toàn $0,25\,\text{kg}$ ở nhiệt độ sôi cần bao nhiêu kJ? Chỉ ghi số.
\shortans{500}
\loigiai{
$Q=Lm=2,0.10^6.0,25=5,0.10^5 J=500 kJ.$
}
\end{ex}

% ===== Câu 94 =====
% ID: L12C1B4-130-TL
% Mức: VDC
\begin{bt}
% ID: L12C1B4-130-TL
% Mức: VDC
Giải thích vì sao khi nước đang sôi ở áp suất không đổi, tiếp tục cấp nhiệt nhưng nhiệt độ gần như không tăng.
\loigiai{
Nhiệt lượng cung cấp chủ yếu làm tăng thế năng tương tác giữa các phân tử để chuyển nước từ lỏng sang hơi, không làm tăng đáng kể động năng trung bình; vì vậy nhiệt độ gần như không đổi.
}
\end{bt}

% ===== Câu 95 =====
% ID: L12C1B4-131-TLN
% Mức: VDC
\begin{ex}
% ID: L12C1B4-131-TLN
% Mức: VDC
Hóa hơi $150\,\text{g}$ chất lỏng cần 135 kJ. Tính $L$ theo đơn vị 10^ $5\,\text{J}$ /kg. Chỉ ghi số.
\shortans{9}
\loigiai{
$m=0,150 kg; L=135000/0,150=9,0.10^5 J/kg.$
}
\end{ex}

% ===== Câu 126 =====
% ID: L12C1B4-132-TN
% Mức: NB
\begin{ex}
% ID: L12C1B4-132-TN
% Mức: NB
Nhiệt hóa hơi riêng của một chất lỏng là nhiệt lượng cần để hóa hơi hoàn toàn bao nhiêu kilôgam chất đó ở nhiệt độ sôi?
\choice
{\True $1\,\text{kg}$}
{$2\,\text{kg}$}
{$1\,\text{g}$}
{1 mol}
\loigiai{
Theo định nghĩa, nhiệt hóa hơi riêng $L$ là nhiệt lượng cần để hóa hơi hoàn toàn $1\,\text{kg}$ chất lỏng ở nhiệt độ sôi.
}
\end{ex}

% ===== Câu 129 =====
% ID: L12C1B4-133-TN
% Mức: NB
\begin{ex}
% ID: L12C1B4-133-TN
% Mức: NB
Trong quá trình sôi ở áp suất không đổi, nhiệt độ của chất lỏng tinh khiết
\choice
{tăng đều}
{giảm đều}
{\True gần như không đổi}
{luôn bằng 0 ° $C$}
\loigiai{
Nhiệt lượng cung cấp dùng cho chuyển thể nên nhiệt độ gần như không đổi.
}
\end{ex}

% ===== Câu 132 =====
% ID: L12C1B4-134-TN
% Mức: NB
\begin{ex}
% ID: L12C1B4-134-TN
% Mức: NB
Công thức Q=Lm chỉ áp dụng trực tiếp khi
\choice
{chất rắn đang nóng chảy}
{\True chất lỏng ở nhiệt độ sôi chuyển hoàn toàn thành hơi}
{chất lỏng chỉ tăng nhiệt độ}
{hơi đang tăng nhiệt độ}
\loigiai{
Q=Lm mô tả riêng phần chuyển thể lỏng thành hơi ở nhiệt độ sôi.
}
\end{ex}

% ===== Câu 135 =====
% ID: L12C1B4-135-TN
% Mức: TH
\begin{ex}
% ID: L12C1B4-135-TN
% Mức: TH
Đơn vị của nhiệt hóa hơi riêng trong hệ SI là
\choice
{$J$}
{\True J/kg}
{J/(kg.K)}
{W/kg}
\loigiai{
Từ L=Q/m, đơn vị của $L$ là J/kg.
}
\end{ex}

% ===== Câu 138 =====
% ID: L12C1B4-136-TN
% Mức: TH
\begin{ex}
% ID: L12C1B4-136-TN
% Mức: TH
Khi hơi của một chất ngưng tụ hoàn toàn ở nhiệt độ không đổi, độ lớn nhiệt lượng tỏa ra là
\choice
{\True Lm}
{mcΔt}
{L/m}
{0}
\loigiai{
Quá trình ngưng tụ tỏa nhiệt có độ lớn Q=Lm.
}
\end{ex}

% ===== Câu 139 =====
% ID: L12C1B4-137-TN
% Mức: TH
\begin{ex}
% ID: L12C1B4-137-TN
% Mức: TH
Nếu cùng một chất ở cùng điều kiện, khối lượng hóa hơi tăng gấp đôi thì nhiệt lượng cần thiết
\choice
{giảm một nửa}
{không đổi}
{\True tăng gấp đôi}
{tăng gấp bốn}
\loigiai{
Do Q=Lm nên $Q$ tỉ lệ thuận với m.
}
\end{ex}

% ===== Câu 140 =====
% ID: L12C1B4-138-TN
% Mức: VD
\begin{ex}
% ID: L12C1B4-138-TN
% Mức: VD
Hệ thức tính nhiệt lượng hóa hơi hoàn toàn khối lượng m ở nhiệt độ sôi là
\choice
{Q=mcΔt}
{$Q=L/m$}
{\True $Q=Lm$}
{$Q=m/L$}
\loigiai{
Khi chất lỏng đã ở nhiệt độ sôi và hóa hơi hoàn toàn: Q=Lm.
}
\end{ex}

% ===== Câu 141 =====
% ID: L12C1B4-139-TN
% Mức: VD
\begin{ex}
% ID: L12C1B4-139-TN
% Mức: VD
Đại lượng nào phụ thuộc bản chất chất lỏng và điều kiện áp suất?
\choice
{Khối lượng m}
{\True Nhiệt hóa hơi riêng $L$}
{Thời gian t}
{Công suất $P$}
\loigiai{
$L$ là đặc trưng của chất và có thể thay đổi theo điều kiện, đặc biệt là áp suất.
}
\end{ex}

% ===== Câu 142 =====
% ID: L12C1B4-140-TN
% Mức: VDC
\begin{ex}
% ID: L12C1B4-140-TN
% Mức: VDC
Nước có L=2,3.10^ $6\,\text{J}$ /kg. Ý nghĩa đúng là
\choice
{$1\,\text{kg}$ nước tăng 1 $K$ cần 2,3.10^6 $J$}
{\True $1\,\text{kg}$ nước ở nhiệt độ sôi hóa hơi hoàn toàn cần 2,3.10^6 $J$}
{$1\,\text{kg}$ hơi nước ngưng tụ luôn thu 2,3.10^6 $J$}
{$1\,\text{kg}$ nước nóng chảy cần 2,3.10^6 $J$}
\loigiai{
Giá trị $L$ cho biết nhiệt lượng hóa hơi của $1\,\text{kg}$ chất lỏng ở nhiệt độ sôi.
}
\end{ex}

% ===== Câu 143 =====
% ID: L12C1B4-141-TN
% Mức: VDC
\begin{ex}
% ID: L12C1B4-141-TN
% Mức: VDC
So với nhiệt lượng làm $1\,\text{kg}$ nước tăng 1 $K$, nhiệt lượng hóa hơi $1\,\text{kg}$ nước ở 100 ° $C$
\choice
{nhỏ hơn rất nhiều}
{bằng nhau}
{\True lớn hơn rất nhiều}
{luôn bằng 0}
\loigiai{
Nhiệt hóa hơi riêng của nước khoảng 2,3.10^ $6\,\text{J}$ /kg, lớn hơn nhiều so với c.1 K.
}
\end{ex}
